\DocumentMetadata{
  pdfversion=2.0,
  pdfstandard=ua-2,
  lang=en-US,
  tagging=on,
  tagging-setup={math/setup=mathml-SE,tabsorder=structure}
}
\documentclass[11pt,letterpaper]{article}
\usepackage[margin=0.68in,includefoot]{geometry}
\usepackage{newcomputermodern}
\usepackage{amsmath,amscd,mathtools}
\providecommand{\square}{\mdlgwhtsquare}
\usepackage[hidelinks]{hyperref}
\hypersetup{
  pdftitle={Numerical Analysis PhD Exam, August 2016, Part 2},
  pdfauthor={University of Florida Department of Mathematics},
  pdfsubject={Graduate mathematics examination},
  pdfdisplaydoctitle=true,
  bookmarksopen=true
}
\setcounter{secnumdepth}{0}
\setlength{\parindent}{0pt}
\setlength{\parskip}{0.28em}
\setlength{\emergencystretch}{3em}
\setlength{\leftmargini}{2.15em}
\setlength{\leftmarginii}{2.65em}
\setlength{\leftmarginiii}{3em}
\renewcommand{\labelenumi}{\textbf{\theenumi.}}
\pagestyle{empty}
\raggedbottom
\allowdisplaybreaks
\newenvironment{examproblems}[1]{%
  \begin{enumerate}%
  \setcounter{enumi}{#1}%
  \setlength{\topsep}{0.35em}%
  \setlength{\partopsep}{0pt}%
  \setlength{\itemsep}{0.48em}%
  \setlength{\parsep}{0.18em}%
}{\end{enumerate}}
\newenvironment{examsubparts}{%
  \begin{enumerate}%
  \setlength{\topsep}{0.18em}%
  \setlength{\partopsep}{0pt}%
  \setlength{\itemsep}{0.16em}%
  \setlength{\parsep}{0.1em}%
}{\end{enumerate}}
\newenvironment{examinstructions}{%
  \par\begingroup\itshape
}{%
  \par\endgroup\vspace{0.18em}
}
\makeatletter
\renewcommand\section{\@startsection{section}{1}{0pt}%
  {-1.25ex plus -.25ex minus -.1ex}%
  {0.55ex plus .1ex}%
  {\normalfont\large\bfseries}}
\renewcommand{\@maketitle}{%
  \newpage\null\vskip -1.2em%
  {\footnotesize\bfseries
    \href{https://www.ufl.edu/}{University of Florida}, \href{https://math.ufl.edu/}{Department of Mathematics}\par}%
  \vskip 0.18em%
  \tagstructbegin{tag=Title}%
    {\LARGE\bfseries\href{https://gma.math.ufl.edu/exams/numerical-analysis-phd/}{Numerical Analysis PhD Exam}, August 2016, Part 2\par}%
  \tagstructend%
  \vskip 0.35em%
  \tagmcbegin{artifact}%
    \hrule height 1pt%
  \tagmcend%
  \vskip 0.55em%
}
\makeatother
\title{Numerical Analysis PhD Exam, August 2016, Part 2}
\author{}
\date{}
\begin{document}
\maketitle
\thispagestyle{empty}
\begin{examinstructions}
You have 2 hours to answer the following questions. You must show all your work as neatly and clearly as possible and indicate the final answer clearly. You may not use a calculator.
\end{examinstructions}
\begin{examproblems}{0}
\item
The iteration 
\[
x_{n+1}=x_n-\frac{f(x_n)}{f'(x_n)-\frac{1}{2}f''(x_n)\frac{f(x_n)}{f'(x_n)}}
\]
 for solving the equation \(f(x)=0\) is known as Halley’s method.
\begin{examsubparts}
\item[\textnormal{(a)}]
Show that the method can, alternatively, be interpreted as applying Newton’s method to the equation \(g(x)=0\), with \(g(x)=f(x)/\sqrt{f'(x)}\).
\item[\textnormal{(b)}]
Assuming~\(\alpha\) is a simple root of the equation, and \(x_n\to\alpha\) as \(n\to\infty\), show that convergence is at least cubic. Hint: Part (a) might help.
\end{examsubparts}
\item
Do the following parts which are unrelated.
\begin{examsubparts}
\item[\textnormal{(a)}]
Derive a Taylor method of order two for the first order initial value problem (IVP) 
\[
\begin{cases}
y'=\dfrac{y^2}{1+t},\\
y(1)=-\dfrac{1}{\ln 2}.
\end{cases}
\tag{1}
\]
\item[\textnormal{(b)}]
Derive a numerical method for the following second order IVP by replacing the derivatives with centered differences over the mesh \(t_{n+1}=t_n+h\) where \(h=1/N\). 
\[
\begin{cases}
y''+2y'+y=\cos t, & 0\le t\le 1,\\
y(0)=1,\\
y'(0)=0.
\end{cases}
\tag{2}
\]
\end{examsubparts}
\item
Consider a quadrature rule of the form 
\[
\int_0^1 f(x)\,dx\approx Af(0)+Bf'(0)+Cf(\gamma)+Df(1).
\]
\begin{examsubparts}
\item[\textnormal{(a)}]
Determine the constants \(A,B,C,D\), and~\(\gamma\) so that the quadrature formula has maximum degree of exactness.
\item[\textnormal{(b)}]
What is the degree of exactness of the formula in part (a)?
\end{examsubparts}
\item
For the function \(f(x)=\ln(x)\) for \(x\in[1,2]\), find the minimax approximation polynomial of degree one. Give the exact value of the minimax error.
\item
This problem has the following parts.
\begin{examsubparts}
\item[\textnormal{(a)}]
Relative to the \(L^2\) norm on the interval \([1,3]\), find the least squares approximation to the function \(f(x)=1/x\) using polynomials of degree at most one.
\item[\textnormal{(b)}]
Find the least squares error of the approximation above.
\end{examsubparts}
\end{examproblems}
\end{document}
